\documentclass[10pt]{article}

\usepackage[a4paper,margin=0.72in]{geometry}
\usepackage{amsmath,amssymb,bm}
\usepackage{booktabs}
\usepackage{array}
\usepackage{enumitem}
\usepackage{xcolor}
\usepackage{microtype}
\usepackage{hyperref}
\usepackage{mathtools}

\hypersetup{colorlinks=true,linkcolor=black,urlcolor=blue,citecolor=black}
\setlength{\parindent}{0pt}
\setlength{\parskip}{4pt}
\setlist[itemize]{leftmargin=1.4em,itemsep=2pt,topsep=2pt}

\newcommand{\Gc}{\mathcal{G}^{c}}
\newcommand{\Da}{\mathcal{D}^{a}}
\newcommand{\Anch}{\mathcal{A}}
\newcommand{\Win}{\mathcal{A}_{W}(t)}
\newcommand{\NKi}{\mathcal{N}_{i}^{(K)}}
\newcommand{\R}{\mathbb{R}}

\title{\vspace{-1.0em}\textbf{STAG-GS:
Sparse Temporal Anchor Learning for Dynamic 3D Gaussian Splatting}\\
\large Sparse State Storage and Query-Time Motion Interpolation}
\author{}
\date{}

\begin{document}
\maketitle
\vspace{-1.2em}

\section*{1. Core Idea}

The framework represents a dynamic scene using \textbf{one persistent 3D Gaussian Splatting (3DGS) scene} together with \textbf{sparse temporal motion records}. Instead of storing a complete set of Gaussians for every frame, only significant changes are retained at selected temporal anchors.

At a query time $t$, the required Gaussian state is reconstructed \emph{on demand}. The method uses:
\begin{itemize}
    \item an \textbf{anchor stride} $A$ controlling temporal storage density,
    \item a \textbf{temporal window} $W$ controlling how many neighboring anchor intervals contribute, and
    \item a \textbf{spatial neighborhood size} $K$ controlling how many nearby Gaussians provide local motion evidence.
\end{itemize}

The output is the interpolated position, rotation, and scale
\[
\hat{\mathbf{x}}_i(t),\qquad \hat{\mathbf{r}}_i(t),\qquad \hat{\mathbf{s}}_i(t),
\]
which are passed directly to the standard 3DGS projection and rasterization pipeline.

\section*{2. Persistent Scene Representation and Sparse Temporal Storage}

For a scene with $N$ persistent Gaussian identities, keep one canonical representation
\[
\Gc = \{G_i^c\}_{i=1}^{N},
\qquad
G_i^c = \bigl(\mathbf{x}_i^c,\mathbf{r}_i^c,\mathbf{s}_i^c,\alpha_i,\mathbf{c}_i,\ldots\bigr),
\]
where $\mathbf{x}_i^c$ is the canonical position, $\mathbf{r}_i^c$ the canonical rotation, $\mathbf{s}_i^c$ the canonical scale, and $\alpha_i,\mathbf{c}_i$ denote the usual 3DGS appearance parameters.

At a retained time $a$, the framework does \textbf{not} save another complete Gaussian scene. Instead, it stores only Gaussians whose motion state differs meaningfully from the canonical state:
\[
\Da =
\left\{
\left(i,\Delta\mathbf{x}_i^a,\Delta\boldsymbol{\rho}_i^a,\Delta\boldsymbol{\ell}_i^a\right)
\;:\;
 d_i^a > \varepsilon
\right\},
\]
with
\begin{align}
\Delta\mathbf{x}_i^a &= \mathbf{x}_i^a-\mathbf{x}_i^c,\\
\Delta\boldsymbol{\rho}_i^a
&= \operatorname{Log}\!\left(\mathbf{r}_i^a\otimes(\mathbf{r}_i^c)^{-1}\right),\\
\Delta\boldsymbol{\ell}_i^a
&= \log \mathbf{s}_i^a-\log \mathbf{s}_i^c.
\end{align}
Here $d_i^a$ is a normalized change score and $\varepsilon$ determines whether the change is significant enough to store. If Gaussian $i$ is absent from $\Da$, its change at anchor $a$ is treated as negligible.

The retained anchor times are
\[
\Anch=\{1,1+A,1+2A,\ldots\}.
\]
Thus, $A=1$ is the dense special case in which every observed frame is eligible as an anchor, while $A>1$ reduces temporal storage.

A rough storage cost is
\[
\underbrace{O(N)}_{\text{canonical 3DGS}}
+
\underbrace{O\!\left(\sum_{a\in\Anch}|\Da|\right)}_{\text{sparse temporal records}},
\]
instead of storing $O(NT)$ complete Gaussian states for $T$ frames. If only a fraction $\rho$ of Gaussians changes significantly at each retained anchor, the dynamic component is approximately
\[
O\!\left(\frac{\rho NT}{A}\right).
\]

\section*{3. Query-Time Context: $A$, $W$, and $K$}

Let query time $t$ lie between two consecutive anchors $a_m\le t\le a_{m+1}$. The temporal window contains neighboring anchors around the query interval:
\[
\Win
=
\{a_{m-W+1},\ldots,a_m,a_{m+1},\ldots,a_{m+W}\}\cap\Anch.
\]
The value $W$ therefore counts \textbf{anchor intervals}, not raw frames. With $W=1$, only the immediate bracketing anchors are used. Increasing $W$ provides more past and future temporal evidence. When $A=1$, every frame is an anchor and $W$ becomes ordinary frame-level temporal context.

For Gaussian $i$, define $\NKi$ as its $K$ nearest Gaussian neighbors in 3D position space. The three main controls therefore have distinct roles:
\[
\boxed{
A=\text{temporal sampling density},\qquad
W=\text{temporal context},\qquad
K=\text{spatial context}.
}
\]

\section*{4. Reconstructing Anchor States Only When Needed}

At inference time, no 3DGS retraining is performed. For each required anchor $a\in\Win$, reconstruct Gaussian $i$ from the canonical state and the sparse record:
\begin{align}
\mathbf{x}_i^a &= \mathbf{x}_i^c+\Delta\mathbf{x}_i^a,\\
\mathbf{r}_i^a &= \operatorname{Exp}(\Delta\boldsymbol{\rho}_i^a)\otimes\mathbf{r}_i^c,\\
\log\mathbf{s}_i^a &= \log\mathbf{s}_i^c+\Delta\boldsymbol{\ell}_i^a.
\end{align}
If no record exists for Gaussian $i$ at anchor $a$, its stored change is taken as zero. Appearance parameters remain in the persistent 3DGS representation; the present formulation focuses on geometric motion through position, rotation, and scale.

\section*{5. Motion State Between Neighboring Anchors}

For two consecutive retained times $a<b$, with $\Delta t=b-a$, define the Gaussian motion rates
\begin{align}
\mathbf{v}_i &= \frac{\mathbf{x}_i^b-\mathbf{x}_i^a}{\Delta t},\\
\boldsymbol{\omega}_i
&= \frac{\operatorname{Log}\!\left(\mathbf{r}_i^b\otimes(\mathbf{r}_i^a)^{-1}\right)}{\Delta t},\\
\boldsymbol{\eta}_i
&= \frac{\log\mathbf{s}_i^b-\log\mathbf{s}_i^a}{\Delta t}.
\end{align}
Stack them into a unified local motion state
\[
\mathbf{u}_i=
\begin{bmatrix}
\mathbf{v}_i\\
\boldsymbol{\omega}_i\\
\boldsymbol{\eta}_i
\end{bmatrix}
\in\R^9.
\]
This allows translation, rotation, and scale to be handled jointly by the same local motion reconstruction stage.

\section*{6. Local Spatial Motion Gradient}

Nearby Gaussians provide spatial evidence about how the motion state changes around Gaussian $i$. For each neighbor $j\in\NKi$, define
\[
\Delta\mathbf{x}_{ij}=\mathbf{x}_j-\mathbf{x}_i,
\qquad
\Delta\mathbf{u}_{ij}=\mathbf{u}_j-\mathbf{u}_i.
\]
Assume that local motion varies approximately linearly over a sufficiently small neighborhood:
\[
\Delta\mathbf{u}_{ij}\approx \mathbf{J}_i\Delta\mathbf{x}_{ij},
\qquad
\mathbf{J}_i\in\R^{9\times3}.
\]
The local motion gradient is estimated by weighted least squares:
\[
\mathbf{J}_i=
\left(
\sum_{j\in\NKi}
 w_{ij}\,\Delta\mathbf{u}_{ij}\Delta\mathbf{x}_{ij}^{\top}
\right)
\left(
\sum_{j\in\NKi}
 w_{ij}\,\Delta\mathbf{x}_{ij}\Delta\mathbf{x}_{ij}^{\top}
+
\epsilon_J\mathbf{I}
\right)^{-1}.
\]
The matrix $\mathbf{J}_i$ describes how the local 9D motion state varies with 3D spatial displacement. It lets nearby Gaussian motion contribute to a locally corrected estimate of the motion of Gaussian $i$.

\section*{7. Bidirectional Temporal Transport and Consensus}

Let the locally corrected motion estimate be
\[
\tilde{\mathbf{u}}_i=
\begin{bmatrix}
\tilde{\mathbf{v}}_i\\
\tilde{\boldsymbol{\omega}}_i\\
\tilde{\boldsymbol{\eta}}_i
\end{bmatrix}.
\]
A transport step of duration $\delta t$ is
\begin{align}
\mathbf{x}_i' &= \mathbf{x}_i+\delta t\,\tilde{\mathbf{v}}_i,\\
\mathbf{r}_i' &= \operatorname{Exp}(\delta t\,\tilde{\boldsymbol{\omega}}_i)\otimes\mathbf{r}_i,\\
\log\mathbf{s}_i' &= \log\mathbf{s}_i+\delta t\,\tilde{\boldsymbol{\eta}}_i.
\end{align}
Within the selected window $\Win$, information is transported from past anchors toward $t$ and independently from future anchors toward $t$. This produces a forward estimate
\[
(\mathbf{x}_i^F,\mathbf{r}_i^F,\mathbf{s}_i^F)
\]
and a backward estimate
\[
(\mathbf{x}_i^B,\mathbf{r}_i^B,\mathbf{s}_i^B).
\]

A reliability weight $\lambda_i\in[0,1]$, derived from the local motion fit, fuses the two estimates:
\begin{align}
\hat{\mathbf{x}}_i(t)
&=\lambda_i\mathbf{x}_i^F+(1-\lambda_i)\mathbf{x}_i^B,\\
\hat{\mathbf{r}}_i(t)
&=\operatorname{SLERP}\!\left(\mathbf{r}_i^B,\mathbf{r}_i^F,\lambda_i\right),\\
\log\hat{\mathbf{s}}_i(t)
&=\lambda_i\log\mathbf{s}_i^F+(1-\lambda_i)\log\mathbf{s}_i^B.
\end{align}
The reconstructed Gaussian set $\hat{\mathcal{G}}(t)$ is then passed directly to the standard 3DGS rasterizer.

\section*{8. Inference Pipeline}

The complete query-time flow is
\[
\boxed{
\text{query }t
\rightarrow
\text{read anchors in }\Win
\rightarrow
\text{KNN of size }K
\rightarrow
\mathbf{u}_i,\mathbf{J}_i
\rightarrow
\text{bidirectional transport}
\rightarrow
\text{consensus}
\rightarrow
\hat{\mathcal{G}}(t)
\rightarrow
\text{3DGS rasterizer}
}.
\]

The framework is therefore \textbf{lazy-evaluated}: dense temporal Gaussian states are not permanently stored. Motion information is reconstructed only for the anchors and Gaussians required by the current query.

\section*{9. Difference from Common Dynamic 3DGS Representations}

\renewcommand{\arraystretch}{1.15}
\begin{center}
\small
\begin{tabular}{>{\raggedright\arraybackslash}p{0.20\linewidth}
                >{\raggedright\arraybackslash}p{0.36\linewidth}
                >{\raggedright\arraybackslash}p{0.36\linewidth}}
\toprule
\textbf{Approach} & \textbf{Stored temporal representation} & \textbf{Query-time motion}\\
\midrule
Learned deformation
& Canonical 3DGS plus a learned deformation network
& Evaluate a learned function such as $F_{\theta}(\mathbf{x},t)$\\
\addlinespace
Explicit trajectory
& Canonical 3DGS plus per-Gaussian temporal coefficients
& Evaluate a stored polynomial, Fourier, or related trajectory model\\
\addlinespace
\textbf{Lazy-evaluated framework}
& Canonical 3DGS plus sparse significant changes at periodic anchors
& Reconstruct local motion on demand using $A$, $W$, $K$, $\mathbf{J}_i$, and bidirectional consensus\\
\bottomrule
\end{tabular}
\end{center}

\section*{10. Main Ablation Axes}

The formulation exposes three direct and interpretable ablations:
\begin{itemize}
    \item \textbf{$A$ --- anchor stride:} tests dense versus sparse temporal storage. Larger $A$ reduces storage but makes interpolation more difficult.
    \item \textbf{$W$ --- temporal window:} tests short versus long temporal context around the query time.
    \item \textbf{$K$ --- spatial neighborhood:} tests how much local Gaussian context is required for reliable motion reconstruction.
\end{itemize}

The dense special case $A=1$ keeps every observed frame eligible as an anchor. Larger values of $A$ progressively move the framework toward sparse temporal storage while preserving the same query-time reconstruction mechanism.

\end{document}